\documentclass[12pt, english]{article}
%%%%%%%%%%%%%%%%%%%%%%%%%%%%%%%%%%%%%%%%%%%%%%%%%%%%%%%%

% margin %%%%%%%%%%%%%%%%%%%%%%
\usepackage{geometry}
 \geometry{
 a4paper,
 left=1in,
 top=0.9in,
 right=1in,
 bottom=2in,
 }
%%%%%%%%%%%%%%%%%%%%%%%%%%%%%%%%%%%%%%%%%%%%%%%%%%%%%%%%

% page no. %%%%%%%%%%%%%%%%%%%%%%
\usepackage{fancyhdr}
\pagestyle{fancy}
\fancyhf{}
\fancyfoot[C]{\textit{\thepage}}
\setlength{\headheight}{65pt}
%%%%%%%%%%%%%%%%%%%%%%%%%%%%%%%%%%%%%%%%%%%%%%%%%%%%%%%%

% idk-needed %%%%%%%%%%%%%%%%%%%%%%
\usepackage{amsmath,amsthm,amssymb,amsfonts, enumitem, color, comment, graphicx, environ, lipsum, accsupp}
%%%%%%%%%%%%%%%%%%%%%%%%%%%%%%%%%%%%%%%%%%%%%%%%%%%%%%%%

% code and color %%%%%%%%%%%%%%%%%%%%%%
\usepackage{spverbatim}
\usepackage{xcolor}
\usepackage{listings}
\definecolor{bluegrey}{RGB}{0, 0, 200}
\lstset{
        framexleftmargin=0pt,
        frame=shadowbox,
        rulesepcolor=\color{gray},
        xleftmargin=30pt,
        xrightmargin=30pt,
        breaklines=true,
        basicstyle=\ttfamily,
        showstringspaces=false,
        }
%%%%%%%%%%%%%%%%%%%%%%%%%%%%%%%%%%%%%%%%%%%%%%%%%%%%%%%%

% heading %%%%%%%%%%%%%%%%%%%%%%
\renewcommand\thesection{\color{blue}\arabic{section}}
\newcommand{\mysection}[2]{\setcounter{section}{#1}\addtocounter{section}{-1}\section{#2}}
%%%%%%%%%%%%%%%%%%%%%%%%%%%%%%%%%%%%%%%%%%%%%%%%%%%%%%%%

% both side indent %%%%%%%%%%%%%%%%%%%%%%
\newenvironment{myindent}
{\par\leftskip25pt\relax\rightskip53.4cm\relax}
{\par\leftskip0cm\relax\rightskip0cm\relax}
%%%%%%%%%%%%%%%%%%%%%%%%%%%%%%%%%%%%%%%%%%%%%%%%%%%%%%%%

% type ^ %%%%%%%%%%%%%%%%%%%%%%
\newcommand{\power}{$^\wedge$}
%%%%%%%%%%%%%%%%%%%%%%%%%%%%%%%%%%%%%%%%%%%%%%%%%%%%%%%%

% document info %%%%%%%%%%%%%%%%%%%%%%
\lhead{Arav Sarma \\ 15193845}  %replace with your name
\rhead{Homework \textcolor{blue}{3} \\ 2025 Speech Processing}
%%%%%%%%%%%%%%%%%%%%%%%%%%%%%%%%%%%%%%%%%%%%%%%%%%%%%%%%

% begin %%%%%%%%%%%%%%%%%%%%%%
\begin{document}
%wwwwwwwwwwwwwwwwwwwwwwwwwwwwwwwwwwwwwwwwwwwwwwwwwwwwwwwww

% section _______________________________________________
\mysection{3}{Procedures}
\vspace{10pt}
% -__-__-__-__-__-__-__-__-__-__-__-__-__-__-__-__-__-__-

% //////////////////////////////////////////////////////
\subsection*{\emph{1.} Sum of whole numbers} 

\begin{lstlisting}
form: "Sum of numbers between two whole numbers"
	integer starting_value
	integer end_value
endform
if end_value < starting_value
	exitScript: "End value cannot be less than starting value."
endif
sum = 0
for i from starting_value to end_value
    sum += i
endfor
writeInfoLine: "The sum of the numbers between ", starting_value, " and ", end_value, " is ", sum, "."
\end{lstlisting}

% \begin{lstlisting}[rulesepcolor=\color{blue}]
% \end{lstlisting}

\vspace{10pt}
% //////////////////////////////////////////////////////

% \\\\\\\\\\\\\\\\\\\\\\\\\\\\\\\\\\\\\\\\\\\\\\\\\\\\\\\\\\
\subsection*{\emph{2.} Product of whole numbers}

\begin{lstlisting}
form: "Product of numbers between two whole numbers"
	integer starting_value
	integer end_value
endform
if end_value < starting_value
	exitScript: "End value cannot be less than starting value."
endif
prod = 1
for i from starting_value to end_value
    prod *= i
endfor
writeInfoLine: "The product of the numbers between ", starting_value, " and ", end_value, " is ", prod, "."
\end{lstlisting}

% \begin{lstlisting}[rulesepcolor=\color{blue}]
% \end{lstlisting}

\newpage
% \\\\\\\\\\\\\\\\\\\\\\\\\\\\\\\\\\\\\\\\\\\\\\\\\\\\\\\\\\

% \\\\\\\\\\\\\\\\\\\\\\\\\\\\\\\\\\\\\\\\\\\\\\\\\\\\\\\\\\
\subsection*{\emph{3.} appendEmptyLine}

\begin{lstlisting}
@appendEmptyLine: "line"
appendInfoLine: "___"
procedure appendEmptyLine: .text$
	appendInfoLine: .text$, newline$
endproc
\end{lstlisting}

% \begin{lstlisting}[rulesepcolor=\color{blue}]
% \end{lstlisting}

\vspace{10pt}
% \\\\\\\\\\\\\\\\\\\\\\\\\\\\\\\\\\\\\\\\\\\\\\\\\\\\\\\\\\

% \\\\\\\\\\\\\\\\\\\\\\\\\\\\\\\\\\\\\\\\\\\\\\\\\\\\\\\\\\
\subsection*{\emph{4.} Old-fashioned mode of address}

\begin{lstlisting}
procedure giveModeOfAddress: .name$, .gender
	if .gender = 1
		writeInfoLine: "Dear Ms. ", .name$
	else 
		writeInfoLine: "Dear Mr. ", .name$
	endif
endproc
form: "Old-fashioned"
	word name
	choice: "gender", 0;
		option: "woman"
		option: "man"
endform
if name$ = ""
	exitScript: "name cannot be blank"
endif
@giveModeOfAddress: name$, gender
\end{lstlisting}

% \begin{lstlisting}[rulesepcolor=\color{blue}]
% \end{lstlisting}

\newpage
% \\\\\\\\\\\\\\\\\\\\\\\\\\\\\\\\\\\\\\\\\\\\\\\\\\\\\\\\\\

% \\\\\\\\\\\\\\\\\\\\\\\\\\\\\\\\\\\\\\\\\\\\\\\\\\\\\\\\\\
\subsection*{\emph{5.} Revised mode of address}

\begin{lstlisting}
procedure giveModeOfAddress: .name$, .honorific
	if self_state_honorific$ <> " o p t i o n a l" & .honorific = 1
		writeInfoLine: "Dear ", self_state_honorific$, " ", .name$
	elsif .honorific = 1
		writeInfoLine: "Dear ", .name$ 
	elsif .honorific = 2
		writeInfoLine: "Dear Captain ", .name$ 
	elsif .honorific = 3
		writeInfoLine: "Dear Ms. ", .name$ 
	elsif .honorific = 4
		writeInfoLine: "Dear Mr.", .name$ 
	endif
endproc
form
	word name
	optionmenu: "gender", 1;
		option: ""
		option: "woman"
		option: "man"
		option: "non-binary"
		option: "other"
	sentence: "self state gender", " o p t i o n a l"
	optionmenu: "honorific", 1;
		option: ""
		option: "Captain"
		option: "Ms."
		option: "Mr."
	sentence: "self_state_honorific", " o p t i o n a l"
endform
if name$ = ""
	exitScript: "name cannot be blank"
endif
@giveModeOfAddress: name$, honorific
\end{lstlisting}
\vspace{5pt}
I wanted to have a text box for self stated values only if the option for that is chosen on the optionmenu (ie. if optionmenu → prefer to self state, open text field inside form), but praat was not accepting if-statements inside the form. So I tried to keep them outside the form, but then praat couldn't open one form after another.
% \begin{lstlisting}[rulesepcolor=\color{blue}]
% \end{lstlisting}

\vspace{10pt}
% \\\\\\\\\\\\\\\\\\\\\\\\\\\\\\\\\\\\\\\\\\\\\\\\\\\\\\\\\\

% ><><><><><><><><><><><><><><><><><><><><><><><><><><><><><
\subsection*{\emph{6.} Participants of your experiment}
\vspace{10pt}

\begin{itemize}
    \item \textbf{General research question:} \par
    Location of phoneme boundary as a function of orthographic context (and exopure/familiarity with orthography of certain language(s))
    \item \textbf{Addition:} \par
    L2 lexical knowledge context
    \item \textbf{Specific research question:} \par
    My idea is to take both orthography and non-native lexical knowledge as a coupled effect on phonological categorisation. Preliminary thoughts on operationalising this question: looking at phonemic boundary between voiced-plain-aspirated, aspirated-fricative continua for L2 English speakers from Assam (or L1 Assamese speaking). Onset stop contrast on L2 variety is through voicing unlike aspiration in L1 English for Assamese (voicing+aspiration for Boro), Orthograhpic <b>, <d>, <g> are percieved as voiced like other Indian Englishes. This perception gets carried over even when exposed to L1 English. Questions: How much closure devoicing / afterburst duration can be blocked off during phoneme categorisation for a voicing / plain percept (how strong would lexical effect be on a word like `ban' / `pan', and also orthographic; eg: when shown <pan> / <phan or fan>)? 
    \item \textbf{Participants} \par
    Participants are going to be tested online using Zoom online meetings, and recruited through author's contacts, and likely by reaching out to Assam/Northeast related online forums like on reddit. The plan is to have non-native English speakers ideally all from Guwahati with same L1, this is to neutralise individual factors like regional variation, exposure to English. A different experiment setup can be by using a qualtrics form along with online meet, so the participants have more autonomy with the test and excess to stimuli on their device while also making it easier to see if they are in a quiet environment.  
\end{itemize}
% ><><><><><><><><><><><><><><><><><><><><><><><><><><><><><

%wwwwwwwwwwwwwwwwwwwwwwwwwwwwwwwwwwwwwwwwwwwwwwwwwwwwwwwww
\end{document}